\subsection{Strategy 2: exact boundary--radial demand propagation}
\label{subsec:exact-demand-guide}

Length alone does not record where a trace must continue. Strategy~2 keeps
that positional information. In local unit vectors
\[
 u=\left(\frac12,\frac{\sqrt3}{2}\right),\qquad
 v=\left(\frac12,-\frac{\sqrt3}{2}\right),
\]
define
\[
 K(a,b,c)=\conv\{0,au,bv,c(u+v)\}.
\]
Let $B_c(a)$ be the largest outgoing demand $b$ for which $K(a,b,c)$ fits
in a closed unit equilateral triangle, and set
\[
 F_c(a)=\min\{B_c(a),1-a\},\qquad G_c(a)=1-F_c(a).
\]
The admissible set is coordinatewise down-closed. Hence $F_c$ is
nonincreasing, $G_c$ is nondecreasing and extensive, and reflection gives
an exact duality. The piecewise formula has a genuine downward jump; no
argument below uses continuity at that jump.

\subsubsection{How an actual row propagates a demand}

Suppose an actual nonsupercritical row has $A_i\ge z$, radial reach
$C_i\ge c_i$, and $A_i+B_i\le1$. Then its realized demand triple gives
\[
 B_i\le F_{c_i}(z).
\]
If the next incident open traces cover their boundary point, the next actual
incoming reach satisfies
\begin{equation}
 A_{i+1}\ge1-B_i\ge G_{c_i}(z).
 \label{eq:actual-propagation}
\end{equation}
Thus each successive use of \eqref{eq:actual-propagation} is an inequality
for an actual row, not merely a formal composition of scalar functions.

The nonlinear selected-$T_+$ part of every capped map is one universal curve.
After normalizing the output by the radial deficit, its increment is
\[
 \psi(x)=1-\sqrt{1-x+x^2}.
\]
This gives strict concavity and both chord bounds without a separate implicit
differentiation in each branch. Once a chain input crosses a low-root
threshold $e(d)$, the map $G_{1-d}$ gives an output at least $1-e(d)$ and all
later extensive maps preserve it. Section~\ref{sec:strategy2-reader} proves
this normal form and the one-hit and two-threshold routing lemmas.

\subsubsection{One signed center model for CE1 and CE2}

The center equations are common. Put
\[
 0<R<1,\qquad W=1-R,\qquad E=\sqrt{1-RW},
\]
\[
 \eta=1-E,\qquad P=E(1-E),\qquad
 k=\eta+\alpha+\delta.
\]
The normalized trace on $e_{0,1}$ and the possible companion trace on
$e_{5,0}$ are
\[
 I_R=\left[\frac{k}{R},W+\delta\right],\qquad
 I_L=\left[\frac{k}{W},R+\alpha\right].
\]
Their signed surpluses are
\[
 \Delta_R=P-\alpha-W\delta,\qquad
 \Delta_L=P-R\alpha-\delta.
\]
The normalization requires $\Delta_R>0$. The center role is CE1 exactly when
$\Delta_L\le0$, and CE2 exactly when $\Delta_L>0$. Thus the second interval
is not a second equation system; it is the positive part of one signed
quantity.

The six center exits are simultaneously
\[
 \begin{aligned}
 d_0^C&=E-\alpha-\delta,&
 d_1^C&=\delta/R,&
 d_2^C&=\delta,\\
 d_3^C&=\min\{\alpha/R,\delta/W\},&
 d_4^C&=\alpha,&
 d_5^C&=\alpha/W.
 \end{aligned}
\]
The complementary row demand is always $c_i=1-d_i^C$. The complete proof,
including the legacy CE1 and CE2 dictionaries, is in
Section~\ref{sec:signed-ce12-reductions}.

\subsubsection{The branch-level contradictions}

For any one-gap branch, repeated use of \eqref{eq:actual-propagation} gives the
same five actual-row chain $1\to2\to3\to4\to5\to0$. Its right-trace input is
$R-\delta$, its supercritical boundary input is $k/R$, and its five radial
deficits are
\[
 \frac\delta R,\quad \delta,\quad
 \min\left\{\frac\alpha R,\frac\delta W\right\},
 \quad\alpha,\quad\frac\alpha W.
\]
The CE1 and CE2 proofs therefore share the full geometric induction, duality
reduction, and terminal diameter comparison. Only the scalar proof of the
middle three-map inequality differs: the CE1 sign uses the universal
selected-$T_+$ chords, whereas the CE2 sign gives a shorter two-threshold
route. Reflection exchanges $R$ with $W$, $\alpha$ with $\delta$, and reverses
the row order.

The two-gap state is genuinely CE2 because it requires $\Delta_L>0$. Its two
endpoint demands still use the same two deficits $\alpha/W$ and $\delta/R$,
but the two gaps share the same supercritical row, so the exact paired
capped-loss theorem is retained.

In the exactly-one-Vd1/Vd2 branch, the preliminary perimeter budget removes
CE1. Every surviving CE2 candidate satisfies
\[
 \alpha+\delta<\frac1{24},\qquad
 \alpha+\delta<\frac{\min\{R,W\}}6.
\]
These two common bounds replace the former outer-ratio and nonadjacent CE2
polynomial comparisons. They give $d_2^C<H/3$ in the adjacent placement and
$d_\tau^C<\min\{A,H\}$ for $\tau=2,3,4$ in the nonadjacent placements.

All exact local-map formulae and the CE1 terminal polynomial estimate remain
in Appendix~\ref{sec:exact-local-demand-calculus}. The signed common center,
short CE2 threshold proof, common budgets, and shortened Vd1/Vd2 placements
are proved in Section~\ref{sec:signed-ce12-reductions}. The half-edge radial
envelope is used only on its valid domain
\[
 p\ge\frac12,\qquad0\le h\le p,\qquad p+h<1;
\]
Subsection~\ref{sec:half-edge-envelope} records its exact counterexample
outside that domain.

\paragraph{Branches closed by exact demand.}
Proposition~\ref{prop:reader-demand-branches} closes every row of
Table~\ref{tab:routing} assigned to route $2$, including the complementary
exact placements in the hybrid $1+2$ row. Singleton gaps are included, and
Equation~\eqref{eq:actual-propagation} is the handoff from every scalar
inequality to the actual role geometry.
